\documentclass[11pt]{article}
\usepackage{titlesec}
\titleformat{\section}[hang]{\normalfont\scshape}{\thesection.}{1em}{}
\titleformat{\subsection}[hang]{\normalfont\scshape}{\thesubsection.}{1em}{}
\usepackage[T1]{fontenc}
\usepackage[dvipsnames]{xcolor}
\usepackage[margin=2cm]{geometry}
\usepackage{amssymb} 
\usepackage{amsmath}
\usepackage{graphicx}
\usepackage{float}
\usepackage[normalem]{ulem} 
\usepackage{enumitem} 
\setlist[enumerate]{itemsep=0mm}
\usepackage{xcolor}
\setenumerate{label=(\arabic*)}
\usepackage[utf8]{inputenc}
\usepackage[T1]{fontenc}
\usepackage{babel}
\usepackage{mathtools}
\usepackage{amsthm}
\usepackage{thmtools}
\usepackage{etoolbox}
\usepackage{fancybox}
\usepackage{framed}
\usepackage{tcolorbox}
\usepackage{xcolor} 
% open 
\newcommand{\open}[0]{\mathcal{O}}
\newcommand{\topo}[0]{\mathcal{T}}
\newcommand{\hood}[0]{\mathcal{U}}
\newcommand{\base}[0]{\mathcal{B}} 
% example environmentq
\theoremstyle{definition} 
\newtheorem{exmp}{Example}[section]

% linear operators
\newcommand{\lop}[2]{\mathcal{L}(#1, #2)}

% question environment
\theoremstyle{definition}
\newtheorem{question}{Question}

%rd 
\newcommand{\rd}[0]{\mathbb{R}^d}
\newcommand{\R}[0]{\mathbb{R}}
\newcommand{\N}[0]{\mathbb{N}} 
% let E in R be measurable 
\newcommand{\EinR}[0]{Let $E \subseteq \R$ be measurable}

%integral 
\newcommand{\idx}[2]{\int_{#1}^{#2}}

% weak convergence 
\newcommand{\warrow}[0]{\rightharpoonup}
\newcommand{\fcvw}[0]{ \{f_n \} \warrow f \text{ in } L^p(E)} 

% Probability 
\DeclareRobustCommand{\bbone}{\text{\usefont{U}{bbold}{m}{n}1}}
\newcommand{\Var}[1]{\mathrm{Var[#1]}}			% variance
\newcommand{\EX}[1]{\mathbb{E}\mathrm{[#1]}}	 % expected value 
\newcommand{\seq}[1]{\{ #1_n	\}_{n \in \bb{N}}} % sequence of events
\newcommand{\pspace}[0]{( \Omega, F, P)}		% probability space
\newcommand{\msp}[0]{( \Omega, F)}		% measurable space
	
% Exercise environment 
\newenvironment{myleftbar}{%
\def\FrameCommand{\hspace{0.6em}\vrule width 2pt\hspace{0.6em}}%
\MakeFramed{\advance\hsize-\width \FrameRestore}}%
{\endMakeFramed}
\declaretheoremstyle[
spaceabove=6pt,
spacebelow=6pt
headfont=\normalfont\bfseries,
headpunct={} ,
headformat={\cornersize*{2pt}\ovalbox{\NAME~\NUMBER\ifstrequal{\NOTE}{}{\relax}{\NOTE}:}},
bodyfont=\normalfont,
]{exobreak}

\declaretheorem[style=exobreak, name=Exercise,%
postheadhook=\leavevmode\myleftbar, %
prefoothook = \endmyleftbar]{exo}

% Solution environment 
\newenvironment{mysolbar}{%
\def\FrameCommand{\hspace{0.6em}\vrule width 2pt\hspace{0.6em}}%
\MakeFramed{\advance\hsize-\width \FrameRestore}}%
{\endMakeFramed}
\declaretheoremstyle[
spaceabove=6pt,
spacebelow=6pt
headfont=\normalfont\bfseries,
headpunct={} ,
headformat={\cornersize*{2pt}\ovalbox{\NAME~\NUMBER\ifstrequal{\NOTE}{}{\relax}{\NOTE}:}},
bodyfont=\normalfont,
]{solbreak}

\declaretheorem[style=solbreak, name=Solution,%
postheadhook=\leavevmode\mysolbar, %
prefoothook = \endmysolbar]{sol}

% HEADERS
\usepackage{fancyhdr}
 
\pagestyle{fancy}
\fancyhf{}
\fancyhead[LE,RO]{Page \thepage}
\fancyhead[RE,LO]{Shreve II}
\fancyhead[CE,CO]{Summer 2026 -- Summary}
\fancyfoot[LE,RO]{}

% Definitions
\newcommand{\dfn}[1]{\underline{\textbf{#1}}}
\newcommand{\im}[1]{\textbf{\textcolor{red}{#1}}}

% lower integral
\usepackage{accents}

\newcommand{\ubar}[1]{\underaccent{\bar}{#1}}
\def\avint{\mathop{\,\rlap{-}\!\!\int}\nolimits} 

% custom commands 
\newcommand{\bb}[1]{\mathbb{#1}}
\newcommand{\vc}[1]{\mathbf{#1}}
\newcommand{\step}[1]{\textbf{#1}\textbf{. Step:}}
\newcommand{\pdv}[2]{\frac{\partial #1}{\partial #2}}
\newcommand{\sets}[2]{ \left\{ #1\ |\ #2 \right\}}
\DeclareMathOperator{\Tr}{Tr}

% Custom operators for Advanced cal
\DeclareMathOperator{\Vol}{Vol}
\DeclareMathOperator{\Content}{Content}
\DeclareMathOperator{\Area}{Area}
\DeclareMathOperator{\curl}{curl}
\DeclareMathOperator{\dive}{div}

% Proofs
\newcommand{\claim}[1]{\textbf{#1}\textbf{. Claim:}}

	% iff proofs
	\newcommand{\rhs}[0]{(\Rightarrow )}
	\newcommand{\lhs}[0]{(\Leftarrow )}

% sequence of functions
\newcommand{\funcseqx}{(f_n(x))_{n \in \bb{N}}}
\newcommand{\funcseq}{(f_n)_{n \in \bb{N}}}

% measurable sets 
\newcommand{\measurable}{f^{-1}([-\infty, c[)} 

% heat equation 
\newcommand{\pbdry}[2]{C^{(#1, #2)} (\Omega_T) \cap C (\overline{\Omega_T})}
\DeclareMathOperator\erf{erf}
\newcommand{\mbf}[1]{\mathbf{#1}}

% Laplace Equation 
\newcommand{\lapbdry}[1]{C^{#1} (\Omega) \cap C (\overline{\Omega})}


% math environments 
\usepackage[utf8]{inputenc}
\newtheorem{theorem}{\textcolor{blue}{Theorem}}
\newtheorem{corollary}{Corollary}
\newtheorem{lemma}[theorem]{Lemma}
\theoremstyle{definition}
\newtheorem{definition}{\textcolor{OliveGreen}{Definition}}
\newtheorem{prop}{\textcolor{red}{Proposition}}
\newtheorem{ex}{\textcolor{Maroon}{Example}}
\theoremstyle{remark}
\newtheorem*{remark}{Remark}

% cookbook proofs 
\newcommand{\cb}[3]{\underline{(#1 #2): #3:}}

\usepackage{tcolorbox}
\tcbuselibrary{theorems}

% theorems 
\newtcbtheorem[number within=section]{mytheo}{Theorem}%
{colback=blue!5,colframe=blue!35!black,fonttitle=\bfseries}{th}

% definitions 
\newtcbtheorem[number within=section]{defn}{Definition}%
{colback=black!5,colframe=black!35!black,fonttitle=\bfseries}{th}

% axioms
\newtcbtheorem[number within=section]{ax}{Axioms}%
{colback=OliveGreen!5,colframe=black!35!OliveGreen,fonttitle=\bfseries}{th}


% important examples
\newtcbtheorem[number within=section]{examp}{Example}%
{colback=Mahogany!5,colframe=black!35!Mahogany,fonttitle=\bfseries}{th}

% upper and lower riemann integrals
\newcommand{\upRiemannint}[2]{
  \overline{\int_{#1}^{#2}}
}
\newcommand{\loRiemannint}[2]{
  \underline{\int_{#1}^{#2}}
}

\usepackage{hyperref}
\hypersetup{
    colorlinks,
    citecolor=black,
    filecolor=black,
    linkcolor=black,
    urlcolor=black
}

\begin{document}

\begin{center}
	\textbf{Shreve: Stochastic Calculus for Finance II} \\
	\textbf{Key Results, Theorems, Definitions, etc.} \\
	\textbf{Charles Zitella}
\end{center}

\begin{abstract}
	This document contains a summary of key results, theorems, and definitions that I deem important or interesting.
    This is not a comprehensive summary of the textbook, but a taste.
\end{abstract}

\tableofcontents

\section{General Probability Theory}

\subsection{Infinite Probability Spaces}

Used to model a situation in which a random experiment with infinitely many possible outcomes is conducted.
The problem with uncountably infinite sample spaces, like uniform sampling from the interval $[0, 1]$, is
that for most experiments the probability of any particular outcome is $0$. 

\begin{definition}[$\sigma$-algebra]
    Let $\Omega$ be a nonempty set, and $\mathcal{F}$ a collection of subsets of $\Omega$. We say that $\mathcal{F}$
    is a \textbf{$\sigma$-algebra} if the following holds:
    \begin{enumerate}
        \item $\emptyset \in \mathcal{F}$ 
        \item $A \in \mathcal{F} \Rightarrow A^c \in \mathcal{F}$ 
        \item $A_1, A_2, \dots \in \mathcal{F} \Rightarrow \bigcup_{n = 1}^\infty A_n \in \mathcal{F}$          
    \end{enumerate}    
    Note that this also implies closed under finite unions and closed under countable intersections (since intersection is complement of union of complements).
    Additionally the whole space $\Omega = \emptyset^c \in \mathcal{F}$.
\end{definition}

\begin{definition}[probability measure]
    Let $\Omega$ nonempty set, $\mathcal{F}$ a $\sigma$-algebra of subsets of $\Omega$. A \textbf{probability measure} $\mathbb{P}$
    is a function that assigns a number in $[0, 1]$ to every set $A \in \mathcal{F}$ called the probability of 
    $A$ and written $\mathbb{P}(A)$. We require that
    \begin{enumerate}
        \item $\mathbb{P}(\Omega) = 1$
        \item (countable additivity) If $A_1, A_2, \dots$ is a sequence of disjoint sets in $\mathcal{F}$, then
        \[
            \mathbb{P}\left(\bigcup_{n = 1}^\infty A_n\right) = \sum_{n = 1}^\infty \mathbb{P}(A_n)
        \]    
    \end{enumerate}
    The triple $(\Omega, \mathcal{F}, \mathbb{P})$ is called a \textbf{probability space}. 
\end{definition}

\begin{ex}{\textit{Uniform Lebesque Measure on $[0, 1]$}}
    
    Define the probability of closed intervals $[a,b]$ to be $\mathbb{P}[a,b] = b - a$, $0 \leq a \leq b \leq 1$.
    Note if $a = b$ then $[a, b] = \{a\}$ and has a probability of $0$. Note this also implies 
    $\mathbb{P}(a, b) = \mathbb{P}[a,b] = b - a$. 

    All sets whose probabilities are determined by the defined formula and the properties of probability measures are exactly those
    belonging to the \textbf{Borel $\sigma$-algebra} $\mathcal{B}([0, 1])$, which is the $\sigma$-algebra generated by the set of open or closed sub-intervals. 
    The sets in this $\sigma$-algebra are called \textbf{Borel sets}.
\end{ex}

\begin{ex}{\textit{Infinite, Independent Coin Toss Space}}

    A coin is tossed infinitely many times and we let $\Omega_{\infty}$ be the set of possible outcomes. The probability
    of a head is $p > 0$ and of tails $q = 1 - p > 0$. We can define a probability measure that gives a probability to every set
    that can be described in terms of finitely many tosses, which in turn determines the probability of sets that are
    not describable with finitely many coin tosses. One can then show that every sequence in $\Omega_\infty$ 
    has probability zero. This brings about the notion of \textit{almost sure}. 
\end{ex}

\begin{definition}[almost sure]
    in the probability space $(\Omega, \mathcal{F}, \mathbb{P})$, if a set 
    $A \in \mathcal{F}$ satisfies $\mathbb{P}(A) = 1$ then it occurs \textbf{almost surely}.  
\end{definition}


\subsection{Random Variables and Distributions}

\begin{definition}
    Let $(\Omega, \mathcal{F}, \mathbb{P})$ be a probability space. A \textbf{random variable} is a real valued function
    $X$ defined on $\Omega$ with the property that for every Borel set $B \in \mathcal{B}(\R)$, the subset of $\Omega$ given by
    \[
    \{X \in B\} = \{\omega \in \Omega : X(\omega) \in B\}
    \] is in the $\sigma$-algebra $\mathcal{F}$.       
\end{definition}

\begin{ex}{\textit{Stock Prices}}
    
    Consider the space of infinite, independent coin tosses from before $(\Omega_\infty, \mathcal{F}_\infty, \mathbb{P})$.
    Define stock prices using the formulas 
    \begin{align*}
S_0(\omega)
    &= 4, \\[1em]
S_1(\omega)
    &=
    \begin{cases}
        8, & \omega_1 = H,\\[4pt]
        2, & \omega_1 = T,
    \end{cases} \\[1em]
S_2(\omega)
    &=
    \begin{cases}
        16, & \omega_1 = \omega_2 = H,\\[4pt]
        4, & \omega_1 \neq \omega_2,\\[4pt]
        1, & \omega_1 = \omega_2 = T,
    \end{cases} \\[1em]
S_{n+1}(\omega)
    &=
    \begin{cases}
        2S_n(\omega), & \omega_{n+1}=H,\\[4pt]
        \dfrac12 S_n(\omega), & \omega_{n+1}=T.
    \end{cases}
\end{align*}
These are all random variables and thus we can compute the probability they take certain values, like
$\mathbb{P}\{S_2 = 4\} = \mathbb{P}(A_{HT} \cup A_{TH}) = 2pq$. 
\end{ex}

\begin{definition}[distribution measure]
    Let $X$ be a random variable on a probability space $(\Omega, \mathcal{F}, \mathbb{P})$. The \textbf{distribution measure}
    of $X$ is the probability measure $\mu_X$ that assigns to each Borel subset $B \subset \R$ the mass
    $\mu_X(B) = \mathbb{P}\{X \in B\}$.
\end{definition}

It's important to note that random variables and distributions are two different concepts. Two different random variables
can have the same distribution and a single random variable can have two different distributions.

\begin{ex} \textit{Distributions}
    
    Consider the uniform measure on $[0, 1]$ defined before, and call it $\mathbb{P}$. Now consider $X(\omega) = \omega$,
    and $Y(\omega) = 1- \omega$. These are two distinct RVs, both defined on $\Omega = [0, 1]$, but they share the same
    distribution:
    \[
        \mu_Y[a,b] = \mathbb{P}\{Y \in [a, b]\} = \mathbb{P}\{\omega : a \leq 1 - \omega \leq b\} = \mathbb{P}[1 - b, 1- a]
        = (1 - a) - (1 - b) = b - a = \mu_X[a,b]
    \]    

    Note we can also define a probability measure on $[0, 1]$ such that the distribution of $X$ and $Y$ are different.
    One such is example is the following:
    \[
        \tilde{\mathbb{P}}[a, b] = \int_a^b 2 \omega \; d\omega = b^2 - a^2, \quad 0 \leq a \leq b \leq 1
    \]   
\end{ex}
Once the distribution measure is defined for all closed intervals $[a, b] \subset \R$, it is determined for 
every Borel subset of $\R$.  

We can also record the distribution of an RV in terms of its \textbf{cumulative distribution function} (cdf), which
is defined as $F(x) = \mathbb{P}\{X \leq x\} = \mu_X(-\infty, x]$, $x \in \R$. The cdf and the distribution measure
essentially give the same amount of information, but we can have more detail in two special cases. Firstly
if there's a nonnegative \textbf{density function} $f(x)$ such that 
\[
    \mu_X[a,b] = \mathbb{P}\{a \leq X \leq b\} = \int_a^b f(x) \; dx \quad -\infty < a \leq b < \infty
\] 
Second if there's a \textbf{probability mass function}. In this case there exists a possibly infinite sequence
$x_1, x_2, \dots$ such that the rv X takes a value in the sequence with probability one. In this case $p_i = \mathbb{P}\{X = x_i\}$ 
The mass assigned to a borel set $B \subset \R$ is 
\[
    \mu_X(B) = \sum_{i : x_i \in B} p_i \quad B \in \mathcal{B}(\R)
\]   

\begin{ex}\textit{Standard Normal Random Variable}

    Let 
    \[
        \varphi(x) = \frac{1}{\sqrt{2 \pi}} e^{\frac{-x^2}{2}}
    \] 
    be the standard normal density, and define the cumulative normal distribution function
    \[
        N(x) = \int_{-\infty}^x \varphi(\zeta) \; d\zeta
    \] 
    Note that this is strictly increasing and continuous, so it has a strictly increasing inverse $N^{-1}(y)$.
    We can construct an rv with a normal distribution on any probability space $(\Omega, \mathcal{F}, \mathbb{P})$
    by starting with a uniformly distributed rv, $Y$ and setting $X = N^{-1}(Y)$. Then
    \begin{align*}
        \mu_X[a, b] &= \mathbb{P}\{\omega \in \Omega : a \leq X(\omega) \leq b\} \\
        &= \mathbb{P}\{\omega \in \Omega : a \leq N^{-1}(Y(\omega)) \leq b\} \\
        &= \mathbb{P}\{\omega \in \Omega : N(a) \leq Y(\omega) \leq N(b)\} \\
        &= N(b) - N(a) \\
        &= \int_a^b \varphi(x) \; dx
    \end{align*}  
    This $\mu_X$ is called the \textbf{standard normal distribution}, and any rv with this distribution regardless
    of probability space is called a standard normal rv. 
\end{ex}

\subsection{Expectations}

We want to define a notion of the "average value" of an rv $X$ on the probability space $(\Omega, \mathcal{F}, \mathbb{P})$.
If $\Omega$ is finite, we can expect the "average" of the rv $X$ to be 
\[
    \mathbb{E}X = \sum_{\omega \in \Omega} X(\omega) \mathbb{P}(\omega)
\]  
If $\Omega$ is countably infinite, we define $\mathbb{E}X$ as the same sum but infinite, but what about continous rvs?
This requires the Lebesque integral, which I will not go into the construction of here (Analysis 3). 

\begin{theorem}[Properties of the Lebesque Integral]
    Let $X$ be a random variable on the probability space $(\Omega, \mathcal{F}, \mathbb{P})$.  
    \begin{enumerate}
        \item If $X$ takes on only finitely many values $y_1, y_2, \dots, y_N$, then 
        \[
            \int_{\Omega} X(\omega) \; \mathbb{P}(\omega) = \sum_{k = 1}^N y_k \mathbb{P}\{X = y_k\}
        \]   
        \item (integrability) The r.v. $X$ is integrable iff $\int_{\Omega} | X(\omega) | \; d\mathbb{P}(\omega) < \infty$
        \item (comparison) If $X \leq Y$ almost surely,\footnote{i.e. $\mathbb{P}\{X \leq Y\} = 1$} and if
        $\int_{\Omega} X(\omega) \; d\mathbb{P}(\omega)$ and $\int_{\Omega} Y(\omega) \; d\mathbb{P}(\omega)$
        are defined, then
        \[
            \int_{\Omega} X(\omega) \; d\mathbb{P}(\omega) \leq \int_{\Omega} Y(\omega) \; d\mathbb{P}(\omega)
        \] 
        If $X = Y$ almost surely and these integrals are defined, then they are equal. 
        \item (linearity) If $\alpha, \beta \in \R_{\geq 0}$, $X, Y$ nonnegative and integrable, then
        \[
            \int_{\Omega} (\alpha X(\omega) + \beta Y(\omega)) \; d\mathbb{P}(\omega) = 
            \alpha \int_{\Omega} X(\omega) \; d\mathbb{P}(\omega) + \beta \int_{\Omega} Y(\omega) \; d\mathbb{P}(\omega)
        \]  
    \end{enumerate}
\end{theorem}

\begin{remark}
    Let $A \in \mathcal{F}$, then
    \[
        \int_A X(\omega) \; d\mathbb{P}(\omega) = \int_{\Omega} \mathbb{I}_A(\omega) X(\omega )\; d\mathbb{P}(\omega)
    \]  
\end{remark}

\begin{definition}[expectation]
    Let $X$ be a random variable on the probability space $(\Omega, \mathcal{F}, \mathbb{P})$. The 
    \textbf{expectation} (or expected value) of $X$ is defined to be
    \[
        \mathbb{E}X = \int_{\Omega} X(\omega) \; d\mathbb{P}(\omega)
    \]  
    This definition makes sense if $X$ is integrable, i.e. $\mathbb{E} |X| < \infty$ or $X \geq 0$ almost surely.  
\end{definition}

Since the expectation is defined in terms of a Lebesque integral, all the previous properties mentioned for the Lebesque
integral also hold for the expectation as well as \textbf{Jensen's inequality}, which states that if
$\varphi$ is a convex, real-valued function defined on $\R$, and if $\mathbb{E} |X| < \infty$, then
$\varphi(\mathbb{E} X) \leq \mathbb{E} \varphi(X)$.

\begin{ex}\textit{Counting Heads}
    
    Consider the infinite coin toss space defined previously, and let $\mathbb{P}$ be the probability measure
    that corresponds to probability $\frac{1}{2}$ for heads on each toss. Let
    \[
        Y_n(\omega) = \begin{cases}
            1 \text{ if } \omega_n = H \\
            0 \text{ if } \omega_n = T
        \end{cases}
    \]  
    Then
    \[
        \mathbb{E} Y_n = 1 \cdot \mathbb{P}\{Y_n = 1\} + 0 \cdot \mathbb{P}\{Y_n = 0\} = \frac{1}{2}
    \] 
\end{ex}

\begin{ex}\textit{Rationals and Irrationals}

    Consider the probability space with $\Omega = [0, 1]$ and $\mathbb{P}$ the lebesque measure.
    Define 
    \[
        X(\omega) = \begin{cases}
            1 \text{ if } \omega \text{ is irrational} \\
            0 \text{ if } \omega \text{ is rational}
        \end{cases}
    \] 
    Then 
    \[
        \mathbb{E} X = 1 \cdot \mathbb{P}\{\omega \in [0, 1] : \omega \text{ is irrational}\} + 0 \cdot \mathbb{P}\{\omega \in [0, 1] : \omega \text{ is rational}\}
    \] 
    Note there are only countably many rational numbers, and the probability of each individual one is $0$, and thus
    by countable additivity $\mathbb{P}\{\omega \in [0, 1] : \omega \text{ is rational}\} = 0$. Therefore the probability
    of the irrationals must be $1$, and subsequently $\mathbb{E} X = 1$.

    It's worth noting that the lebesque integral is essential here, because the integral would be undefined with
    the Riemann definition.
\end{ex}

\begin{definition}[Lebesque measure]
    The \textbf{Lebesque measure} on $\R$, $\mathcal{L}$, is a measure that satisfies countable additivity
    and is defined as $\mathcal{L}[a, b] = b - a$ whenever $a \leq b$.   
\end{definition}

\begin{definition}[almost everywhere]
    We say that a property holds \textbf{almost everywhere} (a.e.) if the set of numbers that fail to satisfy
    the property is of Lebesque measure $0$. 
\end{definition}

\begin{theorem}[Riemann vs. Lebesque integrals]
    Letting $f$ be a bounded function defined on $\R$, and $a < b$ in $\R$,
    \begin{enumerate}
        \item The Riemann integrable $\int_a^b f(x) \; dx$ exists iff $f$ is continuous almost everywhere.
        \item If the Riemann integral is defined, then $f$ is Borel measurable\footnote{The preimage of every Borel set is a Borel set} and the Riemann
        and Lebesque integrals are equal.
    \end{enumerate}    
\end{theorem}

\subsection{Convergence of Integrals}

\begin{definition}[convergence of random variables]
    Let $X_1, X_2, \dots$ be a sequence of random variables, all defined on the same probability space.
    We say $\{X_k\}_{k \geq 0}$ \textbf{converges to $X$ almost surely} if
    \[
    \mathbb{P}\left\{\omega \in \Omega : \lim_{k \to \infty} X_k(\omega) = X(\omega) \right\} = 1
    \]
\end{definition}

\begin{ex}\textit{Strong Law of Large Numbers}

    
    Consider the infinite coin toss space with $\mathbb{P}$ chosen such that the probability of each head is $\frac{1}{2}$.
    Let
    \[
        Y_k = \begin{cases}
            1 \text{ if } \omega_k = H \\
            0 \text{ if } \omega_k = T
        \end{cases}
    \]   
    and
    \[
        H_n = \sum_{k = 1}^n Y_k
    \] 
    Then $H_n$ is the number of heads in the first $n$ tosses, and the strong law of large number asserts that
    \[
        \lim_{n \to \infty} \frac{H_n}{n} = \frac{1}{2} \quad \text{almost surely}
    \]   
    Almost surely because there are some sequences where this isn't true, like all heads, but all those sequences
    have measure $0$.
\end{ex}

\begin{definition}[convergence of measurable functions]
    If $\{f_k\}_{k \geq 1}$ is a sequence of real valued, Borel measurable functions on $\R$, then the sequence
    \textbf{converges to $f$ almost everywhere} if 
    \[
    \mathcal{L}\left\{x \in \R : \lim_{k \to \infty} f_k(x) \neq f(x) \right\} = 0
    \]
\end{definition}
Almost everywhere and almost surely convergence is basically the same thing, just almost surely is used with
r.v.s and probability spaces.

\begin{theorem}[monotone convergence]
    If $\{X_n\}_{n \geq 0}$ is a sequence of almost surely montone increasing random variables that converge almost surely, then
    \[
        \lim_{n \to \infty} \int_{-\infty}^\infty X_n(\omega) \; d\mathbb{P}(\omega) = 
        \int_{-\infty}^\infty \lim_{n \to \infty}  X_n(\omega) \; d\mathbb{P}(\omega)
    \]  
    Therefore $\lim_{n \to \infty} \mathbb{E} X_n = \mathbb{E} X $ where $\lim_{n \to \infty} X_n = X$.  
\end{theorem}

Montone convergence theorem allows us to extend the expectation definition for the discrete case to countably infinite
cases, getting that if $X$ takes on countably many values then $\mathbb{E} X = \sum_{k = 0}^\infty x_k \mathbb{P}\{X = x_k\}$  

\begin{ex}\textit{St Petersburg paradox}

    Same infinite coin toss space, define
    \[
        X(\omega) = \begin{cases}
            2 \text{ if } \omega_1 = H \\
            4 \text{ if } \omega_1 = T, \omega_2 = H \\
            8 \text{ if } \omega_1 = T, \omega_2 = T, \omega_3 = H \\
            \vdots \\
            2^k \text{ if } \omega_1 = \omega_2 = \dots = \omega_{k - 1} = T, \omega_k = H \\
        \end{cases}
    \] 
    
    Then this is finite almost surely, since only non finite case is all tails. Then
    \[
        \mathbb{E} X = \sum_{k \geq 1} 2^k \mathbb{P}\{X = 2^k\} = \sum_{k \geq 1} 2^k \frac{1}{2^k} = \infty
    \] 
\end{ex}

\begin{theorem}[dominated convergence]
    If $\{X_n\}_{n \geq 0}$ is a sequence of random variables and there exists a rv $Y$ such that $\mathbb{E} Y < \infty$
    and $|X_n| \leq Y$ almost surely for all $n$, then we can bring the limit inside the integral, similar to the MCT.   
\end{theorem}

\subsection{Computing Expectations}
Computing expectations using the integral definition requires us to somehow turn it into an integral on $\R$.
This can be done using a density function, but if we don't have one we can use the following. Note that the distribution
measure is a probability measure on $\R$. 

\begin{theorem}
    $X$ a rv on probability space $(\Omega, \mathcal{F}, \mathbb{P})$, and $g$ a Borel measurable function on $\R$.
    Then
    \[
        \mathbb{E} |g(X)| = \int_{\R} |g(x)| \; d\mu_X(x)
    \] 
    and if this quantity is finite then
    \[
        \mathbb{E} g(X) = \int_{\R} g(x) \; d\mu_X(x)
    \] 
\end{theorem}
\begin{proof}
    The proof requires first showing that it is true for indicator functions, and then building up through
    simple functions then non-negative borel measurable functions and finally all Borel measurable functions. \qedhere
\end{proof}

We still can't compute this, but there are several cases we can. If $X$ takes only finitely many values then
$\mu_X$ places a mass of size $p_k = \mathbb{P}\{X = x_k\}$ at each $x_k$ so we get the formula
\[
    \mathbb{E}g(X) = \int_{\R} g(x) \;d\mu_X(x) = \sum_{k = 0}^n g(x_k) p_k
\]     

More commonly we have a nonnegative, Borel measurable density function $f$ on $\R$ such that 
\[
    \mu_X(B) = \int_B f(x) \; dx \quad \text{for every Borel measurable set } B
\]   
and in the case when $f$ is bounded and a.e. continuous, we can compute the Riemann integral and compute the expectation.

\begin{theorem}
    $X$ a random variable on $(\Omega, \mathcal{F}, \mathbb{P})$ and $g$ Borel measurable on $\R$. If $X$ has a density then
    \[
        \mathbb{E} |g(X)| = \int_{-\infty}^\infty |g(x)| f(x) \; dx
    \]     
    and if this is finite then
    \[
        \mathbb{E} g(X) = \int_{-\infty}^\infty g(x) f(x) \; dx
    \]     
\end{theorem}
\begin{proof}
    Once again, show true for indicator functions and build your way up to all Borel measurable functions.
\end{proof}

\subsection{Change of Measure}

\begin{theorem}
    Let $(\Omega, \mathcal{F}, \mathbb{P})$ be a prob. space and $Z$ an almost surely nonnegative random variable
    with $\mathbb{E} Z = 1$. For $A \in \mathcal{F}$, define
    \[
        \tilde{\mathbb{P}}(A) = \int_A Z(\omega) \; d\mathbb{P}(\omega)
    \]     
    Then $\tilde{\mathbb{P}}$ is a probability measure, and if $X$ is a nonnegative random variable then
    \[
        \tilde{\mathbb{E}} X = \mathbb{E}[XZ] 
    \]   
    If $Z$ is almost surely strictly positive then we also have
    \[
        \mathbb{E} X = \tilde{\mathbb{E}}\left[\frac{X}{Z}\right] 
    \]  
\end{theorem}
\begin{remark}
    For a rv with positive and negative values we can apply the theorem to its positive and negative parts provided
    the resulting subtraction does not result in $\infty - \infty$. 
\end{remark}
\begin{proof}[Proof of theorem]
    First
    \[
    \tilde{\mathbb{P}}(\Omega) = \int_{\Omega} Z(\omega) \;d\mathbb{P}(\omega) = \mathbb{E} Z = 1
    \]
    and for disjoint $A_1, A_2, \dots \in \mathcal{F}$, since $B_n = \cup_{k = 1}^n A_k$ is increasing and
    $\lim_{n \to \infty} B_n = B_\infty = \cup_{k = 1}^\infty A_k$ we get by monotone convergence theorem  
    \[
        \tilde{\mathbb{P}}\left(\bigcup_{k = 1}^\infty A_k\right) = \int_{\Omega} \mathbb{I}_{B_\infty}(\omega) Z(\omega) \; d\mathbb{P}(\omega)
        = \lim_{n \to \infty} \int_{\Omega} \mathbb{I}_{B_n}(\omega) Z(\omega) \; d\mathbb{P}(\omega)
        = \sum_{k = 1}^\infty \int_{\Omega} \mathbb{I}_{A_k}(\omega) Z(\omega) \; d\mathbb{P}(\omega) 
        = \sum_{k = 1}^\infty \tilde{\mathbb{P}}(A_k)
    \] 

    Thus it is a probability measure, also for the second part we prove for indicators and move our way up, here is
    the proof for indicators, assume $X = \mathbb{I}_A$ 
    \[
        \tilde{\mathbb{E}} X = \tilde{\mathbb{P}}(A) = \int_{\Omega} \mathbb{I}_A(\omega) Z(\omega) \; d\mathbb{P}(\omega)
        = \mathbb{E}[\mathbb{I}_A Z] = \mathbb{E}[XZ]
    \] 
\end{proof}

\begin{definition}
    Two probability measures on the same space $(\Omega, \mathcal{F})$ are said to be \textbf{equivalent} if they agree
    on which sets have probability zero. 

    The previously defined $\mathbb{P}$ and $\tilde{\mathbb{P}}$ are equivalent.  
\end{definition}

\begin{ex}
    Refer to example $4$, our space is $[0, 1]$, $\mathbb{P}$ is the Lebesque measure, and  
    \[
        \tilde{\mathbb{P}}[a,b] = \int_a^b 2 \omega \; d \omega = b^2 - a^2, \quad 0 \leq a \leq b \leq 1
    \] 
    Then since $\mathbb{P}$ is the lebesque measure, integrating with respect to $\omega$ is the same thing as 
    integrating with respect to $\mathbb{P}(\omega)$, so we get
    \[
        \tilde{\mathbb{P}}[a, b] = \int_a^b 2 \omega \; d\mathbb{P}(\omega)
    \]   
    and since this holds for every closed interval it holds for every Borel set, and thus for $B \in \mathcal{B}$ 
    \[
        \tilde{\mathbb{P}}(B) = \int_B 2 \omega \; d\mathbb{P}(\omega)
    \] 
    Then with $Z(\omega) = 2 \omega$, 
    \[
        \mathbb{E} Z = \int_0^1 2 \omega \; d \omega = 1
    \] 
    So we can apply the previous theorem to any non-negative RV $X(\omega)$ and get
    \[
        \int_0^1 X(\omega) \; d \tilde{\mathbb{P}}(\omega) = \tilde{\mathbb{E}} X = \mathbb{E} [XZ] = \int_0^1 X(\omega) \cdot 2 \omega \; d\mathbb{P}(\omega)
    \] 
    and thus
    \[
        d\tilde{\mathbb{P}}(\omega) = 2 \omega \; d \mathbb{P}(\omega)
    \] 
\end{ex}

\begin{definition}[Radon-Nikodym derivative]
    $(\Omega, \mathcal{F}, \mathbb{P})$ a prob. space and $\tilde{\mathbb{P}}$ another equivalent prob. measure
    on the same space, with an almost surely positive rv $Z$ such that $\tilde{\mathbb{P}}(A) = \int_A Z(\omega) \; d \mathbb{P}(\omega)$.
    Then $Z$ is called the \textbf{Radon-Nikodym derivative} of $\tilde{P}$ with respect to $\mathbb{P}$, and we write
    \[
        Z = \frac{d \tilde{\mathbb{P}}}{d\mathbb{P}}
    \]       
\end{definition}


\begin{theorem}[Radon-Nikodym]
    $\mathbb{P}$ and $\tilde{\mathbb{P}}$ equivalent probability measures on the same space $(\Omega, \mathcal{F})$.
    Then there exists an almost surely positive random variable $Z$ such that $\mathbb{E} Z = 1$ and for every $A \in \mathcal{F}$ 
    \[
        \tilde{\mathbb{P}}(A) = \int_A Z(\omega) \; d \mathbb{P}(A)
    \]  
    
\end{theorem}





\subsection{Exercises}


\section{Information and Conditioning}





\end{document}